% Generated draft: requires mathematical and source review.
\documentclass[11pt]{article}
\usepackage[margin=25mm]{geometry}
\usepackage{fontspec}
\setmainfont{Latin Modern Roman}
\usepackage{amsmath,amssymb,mathtools,unicode-math}
\setmathfont{Latin Modern Math}
\usepackage{enumitem,needspace,xurl,hyperref}
\setlength{\parindent}{0pt}
\setlength{\parskip}{5pt}
\setlength{\emergencystretch}{3em}
\newcommand{\N}{\mathbb{N}}
\newcommand{\Z}{\mathbb{Z}}
\newcommand{\Q}{\mathbb{Q}}
\newcommand{\R}{\mathbb{R}}
\newcommand{\C}{\mathbb{C}}
\newcommand{\F}{\mathbb{F}}
\newcommand{\A}{\mathbb{A}}
\newcommand{\PP}{\mathbb{P}}
\newcommand{\E}{\mathbb{E}}
\newcommand{\eps}{\varepsilon}
\newcommand{\dd}{\,\mathrm{d}}
\begin{document}
\section*{UnsolvedMath ID: 800374}
\textbf{Source problem code:} catalogue entry
\subsection{Definitions and mathematical statement}
Let \(d\geq 0\) be an integer and let
\[
H(t)=\sum_{i=0}^{d}h_i t^i
\]
be a polynomial with nonnegative integer coefficients, with \(h_0=1\), whose every complex root is real. Let \(k\) be an integer satisfying \(0\leq k\leq d\) and
\[
h_0<h_1<\cdots<h_k.
\]
For \(k=0\), this chain imposes no additional inequality.

A finite abstract simplicial complex \(\Delta\) is a finite family of finite sets that contains the empty set and is closed under taking subsets. Write
\[
f_j(\Delta)=\#\{F\in\Delta:|F|=j+1\}\qquad(j\geq -1).
\]
In particular, \(f_{-1}(\Delta)=1\). Using the augmented convention, its \(f\)-vector includes this empty-face count.

\textbf{Question (Brenti, 2024, Problem 12).} For every polynomial \(H\) and every integer \(k\) satisfying the preceding conditions, does there exist a finite abstract simplicial complex \(\Delta\), having no face of cardinality greater than \(k\), such that
\[
f_{-1}(\Delta)=h_0=1,
\qquad
f_{i-1}(\Delta)=h_i-h_{i-1}\quad(1\leq i\leq k)?
\]
Equivalently, under this convention, is
\[
(h_0,h_1-h_0,\ldots,h_k-h_{k-1})
\]
the augmented \(f\)-vector of a simplicial complex? For \(k=0\), the requested complex is \(\{\varnothing\}\).

\textit{Unresolved source qualification.} The supplied record does not define its \(f\)-vector convention. The formulation above interprets the initial entry \(h_0=1\) as the empty-face count. If the source instead omits the empty-face count, its displayed tuple would prescribe \(f_0,\ldots,f_k\), giving a different question. The intended convention requires source confirmation.
\subsection{Short English statement}
Brenti asks whether the successive differences along any strictly increasing initial segment of a real-rooted polynomial with nonnegative integer coefficients and constant term \(1\) count the faces of a simplicial complex. The statement above includes the empty face as the first \(f\)-vector entry; this convention is not explicitly specified in the supplied record.
\subsection{Sources}
\begin{itemize}
\item Supplied local UnsolvedMath registry record: permanent ID \(800374\), source problem code \texttt{catalogue entry}, title ``Coxeter Groups and Unimodality --- Problem 12,'' in the AMR Open Problem Lists collection. The record attributes the problem to Brenti, with proposed year \(2024\), and labels it \texttt{open}. This is a supplied status label, not an independently verified mathematical or literature assessment.
\item Brenti, \emph{Some Open Problems on Coxeter Groups and Unimodality} (2024), Problem 12. Supplied source URL: \url{https://arxiv.org/abs/2410.09897}. The background records access on \texttt{2026-10-16} and extraction by \texttt{arxiv\_latex}; these are supplied provenance notes, not a claim of access or verification for this draft.
\item Registry provenance: upstream ID \(800374\), upstream revision \texttt{d371ee917ac9647fbb06cd6fc529deec6f3dda82}; record created and updated at \texttt{2026-10-16T14:00:00Z}.
\end{itemize}
\end{document}
